\documentclass[11pt,a4paper]{article}
\usepackage[utf8]{inputenc}
\usepackage[T1]{fontenc}
\usepackage[brazil]{babel}
\usepackage[margin=2.5cm]{geometry}
\usepackage{mathptmx}
\usepackage{enumitem}
\usepackage{booktabs}
\usepackage{tabularx}
\usepackage{array}
\usepackage{hyperref}
\hypersetup{colorlinks=false, pdfborder={0 0 0}}
\usepackage{fancyhdr}

\pagestyle{fancy}
\fancyhf{}
\fancyhead[L]{\small TrackFleet --- Estrutura Analítica dos Recursos}
\fancyhead[R]{\small\thepage}
\renewcommand{\headrulewidth}{0.4pt}

\newcolumntype{L}{>{\raggedright\arraybackslash}X}
\newcolumntype{C}{>{\centering\arraybackslash}X}
\newcolumntype{R}{>{\raggedleft\arraybackslash}X}
\setlist{topsep=3pt}

% Cabeçalho padrão do projeto. No documento consolidado (trabalho_pratico_01/)
% estes três comandos são redefinidos e este mesmo arquivo vira parte do capítulo do domínio.
\newcommand{\cabecalho}[3]{%
  \begin{center}
  {\Large\bfseries DOCUMENTO DE #1 DO PROJETO}\\[2pt]
  {\large\bfseries TrackFleet --- Gestão de Rastreador de Automóveis}\\[4pt]
  Disciplina de Gerenciamento de Projetos $\cdot$ Framework PMBOK\textregistered\ 8\textsuperscript{a} Edição $\cdot$ Domínio: #2\\[2pt]
  Integrantes: Enzo Allebrand e Kerlon Ribeiro (dois GPs) $\cdot$ 4\textsuperscript{o} semestre $\cdot$ Data: #3
  \end{center}
  \vspace{2mm}\hrule\vspace{4mm}}
\newcommand{\parte}[1]{\noindent\textbf{\large #1}\par\vspace{1mm}}
\newcommand{\separador}{\par\vspace{2mm}\hrule\vspace{4mm}}

\begin{document}

\cabecalho{RECURSOS}{Recursos}{17/09}

\parte{Estrutura Analítica dos Recursos (EAR)}
\noindent Organiza os recursos do projeto por categoria, de forma hierárquica, para estimar, adquirir e monitorar --- deixando claro que recurso não é só gente: serviços, licenças e infraestrutura também têm custo, prazo de aquisição e limite de capacidade.

\section{Composição da EAR}

\begin{itemize}[itemsep=4pt]
  \item[\textbf{1}] Recursos do Projeto
  \begin{itemize}[itemsep=3pt]
    \item[\textbf{1.1}] Equipe (recursos humanos)
    \begin{itemize}[itemsep=1pt]
      \item[\textbf{1.1.1}] Gestão em co-gestão: os dois GPs.
      \item[\textbf{1.1.2}] Backend, motor de alertas, relatórios e infraestrutura: Enzo Allebrand.
      \item[\textbf{1.1.3}] Integração com o provedor, frontend, testes e LGPD: Kerlon Ribeiro.
      \item[\textbf{1.1.4}] Apoio do cliente: tempo do gestor de frota da transportadora para entrevistas, validações, treinamento e piloto.
    \end{itemize}
    \item[\textbf{1.2}] Equipamentos (recursos físicos)
    \begin{itemize}[itemsep=1pt]
      \item[\textbf{1.2.1}] Estações de trabalho: notebooks pessoais dos dois GPs.
      \item[\textbf{1.2.2}] Rastreadores GPS já instalados nos veículos da transportadora piloto; não são comprados nem instalados pelo projeto (restrição de orçamento zero para hardware).
    \end{itemize}
    \item[\textbf{1.3}] Software e serviços (recursos virtuais)
    \begin{itemize}[itemsep=1pt]
      \item[\textbf{1.3.1}] API do provedor de rastreamento: credenciais de teste e de produção, limite de requisições e intervalo de envio --- o recurso mais crítico do projeto.
      \item[\textbf{1.3.2}] Hospedagem: um VPS de baixo custo com Docker, PostgreSQL em contêiner e Nginx com certificado TLS gratuito.
      \item[\textbf{1.3.3}] Mapas: Leaflet com tiles do OpenStreetMap, sem custo de licença, sujeitos à política de uso justo.
      \item[\textbf{1.3.4}] E-mail transacional para as notificações de alerta, em plano gratuito.
      \item[\textbf{1.3.5}] Domínio .com.br.
      \item[\textbf{1.3.6}] Ferramentas de trabalho: GitHub, VS Code e Discord, todos gratuitos.
    \end{itemize}
  \end{itemize}
\end{itemize}

\section{Ligação com a EAP e com Finanças}
Cada pacote da EAP consome recursos de um ou mais ramos da EAR: o desenvolvimento (1.3 da EAP) usa os ramos 1.1.2, 1.1.3, 1.2.1 e 1.3.1; a implantação (1.3.5 da EAP) usa o ramo 1.3.2; o piloto (1.4.2 da EAP) usa 1.1.4, 1.2.2 e 1.3.1. Os ramos 1.3.1, 1.3.2, 1.3.4 e 1.3.5 formam o OpEx de R\$ 700 do Documento de Finanças, custeado pela transportadora parceira.

\section{Estratégia de Controle}
O controle mais rigoroso vai para o ramo 1.3, porque é onde os limites são externos: o limite de requisições da API do provedor, a cota do plano gratuito de e-mail e a política de uso justo dos tiles do OpenStreetMap. Essas cotas são conferidas uma vez por mês e antes do piloto, e o custo mensal do ramo é comparado ao OpEx previsto. As cotas da API também orientam decisões de arquitetura --- por exemplo, a frequência das consultas do job de ingestão.

\end{document}
